\section{Three Object Joins: Source and Measurement Trace}\label{sec:object-join-evidence}
The following examples instantiate established adaptive-evaluation and deployment principles in particular serving and agent objects. They are applications of those principles, not claims that the principles originated here. Juarez et al. examined how shifts in browsing habits, location and browser version affect website-fingerprinting attacks; Mathews et al. evaluated nine defenses with stronger attacks and an on-network DynaFlow implementation. Tram\`er et al. and Carlini et al. give broader adaptive and robustness-evaluation guidance \cite{juarez2014critical,mathews2023critical,tramer2020adaptive,carlini2019evaluating}.

\paragraph{Prefix isolation versus semantic-cache processing}
Early Bird studies both a prefix-recovery oracle and a semantic-cache attack. Its average 6\,ms processing delay and approximately 4\% overhead describe de-identification on a GPTCache semantic-cache hit, whose unprotected mean is approximately 0.14\,s \cite[V-A--B]{earlybird2025}. They do not price its $K$-prefix isolation rule. For contrast, CachePrune reports a default serving configuration with QASPER generation quality 0.28 versus 0.30 and time-to-first-token 273 versus 797\,ms for no sharing \cite[Table~2, 5.2--5.6]{wu2026cacheprune}. Its direct-recovery, quality and efficiency results are kept with their source-defined configuration and threat; the two papers' numbers do not form a controlled cross-system comparison.

\paragraph{Discarded reads versus executed writes}
In a replay of reference programs for 97 eligible retail tasks, each repeated across three cover selections, the read-only and isolated-shadow conditions preserve reference database state in all 291 units at each insertion position. In contrast, a same-item self-write inserted before the reference program changes state in 173/291 units: the tool can change order status and append a refund entry, even when the item identifier and price do not change (Table~\ref{tab:tau-cover-bridge}, S9). These are simulator outcomes for the specified APIs and states. They show why discarding a response does not by itself establish that an executed write has no effect; they are not an estimate of mutation prevalence across real services.

\paragraph{Component quality versus planner outcome}
CachePrune's QASPER score above measures generation quality in its serving evaluation, not an acting agent's task completion. A separate GLM experiment obtains new model decisions for 50 previously known retail tasks in each arm. It records 37/50 native successes without cover and 33/50 with a prefix-shadow scheduler, a paired difference of $-8.0$ percentage points with customer-bootstrap interval $[-23.9,6.5]$ over 37 customers (S26). This interval does not resolve a success-rate difference. The measured scheduler has injected response delays and deadline warnings, and omits the full joint framing and terminal flush used in the trace replay. Its result therefore binds utility to this particular live diagnostic; it cannot be assigned to that other defense configuration or interpreted as a WAN latency result.
